% ═══════════════════════════════════════════════════════════════
%  chapters/ch03-architecture.tex
%  Архитектура проекта: три мира, потоки данных, модули
% ═══════════════════════════════════════════════════════════════

\chapter{Архитектура проекта}
\label{ch:architecture}

\epigraph{<<Хорошая архитектура~--- это когда каждый модуль
можно объяснить одним предложением>>.}{}

\section{Введение}
\label{sec:arch-intro}

В \texttt{ch01} мы увидели \emph{общую} картину: что такое
осциллограф, что такое генератор, зачем нужны \emph{два}
STM32. В \texttt{ch02}~--- \emph{прошили} первую программу
и убедились, что \emph{инструменты} работают.

Теперь~--- \textbf{архитектура}. Прежде чем писать
\emph{код}~--- нужно понять, \emph{как} его писать.
Не <<в одном файле>>, не <<как получится>>, а
\textbf{осознанно}.

\begin{ideabox}
\textbf{Архитектура~--- это не <<схема>>, а \emph{решение}.}
Решение о том, \emph{где} провести границы между модулями.
Если границы проведены \emph{правильно}~--- код \emph{легко}
читать, \emph{легко} тестировать, \emph{легко} развивать.
Если \emph{неправильно}~--- всё превращается в кашу.
\end{ideabox}

\subsection{Что вы узнаете}

\begin{itemize}
    \item Из каких \emph{трёх миров} состоит проект.
    \item Как данные \emph{текут} от осциллографа к графику.
    \item Почему \texttt{AppData}~--- \emph{центральный} хаб состояния.
    \item Как \emph{разбить} <<монолитный>> \texttt{main.c}
          на \emph{модули}.
    \item Как \emph{работает} поток чтения осциллографа.
    \item Почему \emph{многопоточность}~--- это \emph{сложно},
          но \emph{необходимо}.
\end{itemize}

\section{Три мира}
\label{sec:arch-three-worlds}

Проект состоит из \textbf{трёх} независимых \emph{миров},
которые \emph{общаются} друг с другом по \textbf{протоколу}.

\begin{figure}[htbp]
\centering
\adjustbox{max width=\textwidth}{%
\begin{tikzpicture}[
    world/.style={
        draw, rounded corners=3mm, thick,
        minimum width=4cm, minimum height=3cm,
        align=center, font=\small
    },
    osc/.style={world, fill=clrHWData!15},
    gen/.style={world, fill=clrLogicData!15},
    pc/.style={world, fill=clrUIData!15},
    arr/.style={-{Latex[length=4mm]}, very thick, color=clrFlowCmd},
    arrdata/.style={-{Latex[length=4mm]}, very thick, color=clrFlowData},
    label/.style={
        font=\small,
        fill=white,
        inner sep=3pt,
        above,
        midway,
        yshift=8pt
    },
    labelbelow/.style={
        font=\small,
        fill=white,
        inner sep=3pt,
        below,
        midway,
        yshift=-8pt
    }
]

% Осциллограф
\node[osc] (osc) at (-5, 0) {\textbf{Осциллограф}\\[2pt]
    NUCLEO-F401RE\\
    ADC + DMA\\
    USB CDC};

% ПК
\node[pc] (pc) at (0, 0) {\textbf{ПК}\\[2pt]
    GTK4 + Cairo\\
    Потоки + Мьютексы\\
    Протокол};

% Генератор
\node[gen] (gen) at (5, 0) {\textbf{Генератор}\\[2pt]
    STM32F103\\
    DAC + PWM\\
    USB CDC};

% Стрелки: осциллограф → ПК
\draw[arrdata, bend left=20] (osc.north east) to
    node[label] {данные (OSC\_DATA)} (pc.north west);

% Стрелки: ПК → осциллограф
\draw[arr, bend left=20] (pc.south west) to
    node[labelbelow] {команды (stream\_on)} (osc.south east);

% Стрелки: ПК → генератор
\draw[arr, bend right=20] (pc.north east) to
    node[label] {команды (set\_wave)} (gen.north west);

% Стрелки: генератор → ПК
\draw[arrdata, bend right=20] (gen.south west) to
    node[labelbelow] {ответы (get\_status)} (pc.south east);

\end{tikzpicture}
}% end adjustbox
\caption{Три мира проекта. \textcolor{clrFlowData}{Данные}
текут от осциллографа к ПК, \textcolor{clrFlowCmd}{команды}~---
от ПК к генератору и обратно.}
\label{fig:arch-three-worlds}
\end{figure}

\subsection{Мир 1: осциллограф}

\textbf{Осциллограф}~--- на плате \textbf{NUCLEO-F401RE}.
Его \emph{задача}:

\begin{itemize}
    \item \textbf{Захватывать} сигнал через \textbf{ADC}.
    \item \textbf{Перекладывать} данные в память через
          \textbf{DMA} (без участия CPU).
    \item \textbf{Формировать} кадры \texttt{OSC\_DATA}
          с метаданными.
    \item \textbf{Передавать} кадры на ПК через
          \textbf{USB CDC}.
\end{itemize}

\begin{keybox}
\textbf{Осциллограф~--- это <<слушатель>>.} Он \emph{не}
принимает решений, \emph{не} рисует графики. Он \emph{только}
захватывает сигнал и \emph{отдаёт} его ПК. Вся \emph{логика
отображения}~--- на ПК.
\end{keybox}

\subsection{Мир 2: генератор}

\textbf{Генератор}~--- на плате \textbf{STM32F103C8T6}
(Blue Pill). Его \emph{задача}:

\begin{itemize}
    \item \textbf{Принимать} команды от ПК через
          \textbf{USB CDC}.
    \item \textbf{Формировать} сигнал через \textbf{DAC}
          или \textbf{PWM}.
    \item \textbf{Воспроизводить} таблицу формы через
          \textbf{DMA circular}.
    \item \textbf{Отвечать} на запросы статуса.
\end{itemize}

\begin{keybox}
\textbf{Генератор~--- это <<голос>>.} Он \emph{принимает}
команды и \emph{формирует} сигнал. \emph{Не} хранит
историю, \emph{не} рисует графики. Только \emph{выдаёт}
сигнал.
\end{keybox}

\subsection{Мир 3: ПК}

\textbf{ПК}~--- на \textbf{GTK4}. Его \emph{задача}:

\begin{itemize}
    \item \textbf{Принимать} данные от осциллографа
          и \emph{рисовать} их.
    \item \textbf{Отправлять} команды генератору.
    \item \textbf{Управлять} потоком данных (start/stop).
    \item \textbf{Отображать} интерфейс: форму, частоту,
          амплитуду.
\end{itemize}

\begin{keybox}
\textbf{ПК~--- это <<мозг>>.} Он \emph{принимает} решения:
что рисовать, как масштабировать, когда остановить поток.
Вся \emph{логика управления}~--- на ПК.
\end{keybox}

\section{Потоки данных}
\label{sec:arch-data-flow}

Данные в проекте текут в \textbf{двух направлениях}:
\emph{команды}~--- от ПК к устройствам, \emph{данные}~---
от устройств к ПК.

\subsection{Поток команд: сверху вниз}

Пользователь \emph{меняет} форму сигнала в интерфейсе:

\begin{enumerate}
    \item \textbf{GTK4} ловит \emph{сигнал} и вызывает
          \texttt{on\_apply\_generator()}.
    \item \texttt{on\_apply\_generator()} \emph{собирает}
          команды (set\_wave, set\_freq, set\_ampl, ...).
    \item \texttt{send\_cmd()} \emph{формирует} кадры
          и \emph{отправляет} их в порт генератора.
    \item \textbf{STM32F103} \emph{принимает} кадры
          и \emph{применяет} настройки.
    \item \textbf{DAC} \emph{начинает} формировать
          новый сигнал.
\end{enumerate}

\subsection{Поток данных: снизу вверх}

Arduino (осциллограф) \emph{захватывает} сигнал:

\begin{enumerate}
    \item \textbf{ADC} \emph{семплирует} сигнал.
    \item \textbf{DMA} \emph{перекладывает} семплы в SRAM.
    \item \textbf{Прошивка} \emph{формирует} кадр
          \texttt{OSC\_DATA}.
    \item \textbf{USB CDC} \emph{отправляет} кадр на ПК.
    \item \texttt{osc\_reader\_thread()} \emph{принимает}
          кадр в отдельном потоке.
    \item \texttt{draw\_scope()} \emph{рисует} сигнал
          на экране.
\end{enumerate}

\begin{figure}[htbp]
\centering
\adjustbox{max width=\textwidth}{%
\begin{tikzpicture}[
    box/.style={
        draw, rounded corners=1.5mm, thick,
        minimum width=2.4cm, minimum height=0.9cm,
        align=center, font=\footnotesize\ttfamily,
        fill=white
    },
    hw/.style={box, fill=clrHWData!25},
    logic/.style={box, fill=clrLogicData!20},
    ui/.style={box, fill=clrUIData!15},
    arr/.style={-{Latex[length=3mm]}, very thick, color=clrFlowData},
    node distance=1.2cm,
label/.style={
    font=\small,
    fill=white,
    inner sep=3pt,
    above,
    midway,
    yshift=14pt
}
]

% Поток данных (вверх)
\node[hw] (adc) {ADC};
\node[hw, right=of adc] (dma) {DMA};
\node[hw, right=of dma] (frame) {Кадр};
\node[hw, right=of frame] (usb) {USB CDC};
\node[logic, right=of usb] (thread) {Поток};
\node[ui, right=of thread] (draw) {Cairo};

\draw[arr] (adc) -- (dma) node[label] {семплы};
\draw[arr] (dma) -- (frame) node[label] {блок};
\draw[arr] (frame) -- (usb) node[label] {OSC\_DATA};
\draw[arr] (usb) -- (thread) node[label] {read()};
\draw[arr] (thread) -- (draw) node[label] {данные};

\end{tikzpicture}
}% end adjustbox
\caption{Поток данных от ADC до Cairo.}
\label{fig:arch-data-path}
\end{figure}

\begin{keybox}
\textbf{Данные~--- это <<кровь>> проекта.} Команды~--- это
\emph{управление}, данные~--- это \emph{жизнь}. Если поток
данных \emph{остановился}~--- осциллограф \emph{мёртв}.
\end{keybox}

\section{Модули ПК-приложения}
\label{sec:arch-pc-modules}

\textbf{Монолитный} \texttt{main.c} (как в прототипе)~---
это \emph{неправильно}. Мы \emph{разобьём} его на \textbf{шесть}
модулей.

\begin{table}[htbp]
\centering
\caption{Модули ПК-приложения}
\label{tab:arch-pc-modules}
\small
\begin{tabularx}{\textwidth}{@{}l L@{}}
\toprule
\textbf{Модуль} & \textbf{Ответственность} \\
\midrule
\texttt{protocol.c}    & Отправка/приём кадров, CRC16 \\
\texttt{serial\_port.c} & Открытие порта, termios, USB CDC \\
\texttt{osc\_reader.c}  & Поток чтения осциллографа \\
\texttt{draw\_scope.c}  & Отрисовка осциллограммы (Cairo) \\
\texttt{ui.c}           & GTK4-интерфейс, кнопки, поля \\
\texttt{main.c}         & Точка входа, инициализация \\
\bottomrule
\end{tabularx}
\end{table}

\subsection{protocol.c}

\textbf{Что делает}: \emph{формирует} и \emph{парсит}
кадры протокола.

\begin{itemize}
    \item \texttt{send\_cmd()}~--- \emph{отправляет}
          кадр с командой.
    \item \texttt{crc16\_ibm()}~--- \emph{считает} CRC16.
    \item \texttt{parse\_frame()}~--- \emph{парсит}
          входящий кадр.
\end{itemize}

\begin{notebox}
\textbf{Протокол} \emph{подробно} разбираем в
\texttt{ch04}. Здесь~--- только \emph{упоминание}.
\end{notebox}

\subsection{serial\_port.c}

\textbf{Что делает}: \emph{открывает} порт и
\emph{настраивает} \texttt{termios}.

\begin{itemize}
    \item \texttt{open\_port()}~--- \emph{открывает}
          порт с \texttt{O\_NOCTTY}, \texttt{O\_NONBLOCK}.
    \item \texttt{configure\_port()}~--- \emph{настраивает}
          \texttt{termios}.
    \item \texttt{close\_port()}~--- \emph{закрывает} порт.
\end{itemize}

\textbf{Важно}: \emph{тот} \emph{же} \texttt{termios}, что
в \emph{первой} книге. \emph{Принципы}~--- те же, только
\emph{устройство} другое (USB CDC вместо UART).

\subsection{osc\_reader.c: самое интересное}

\textbf{Что делает}: \emph{читает} данные из порта
в \emph{отдельном} потоке.

\begin{listing}[htbp]
\begin{minted}[frame=none, bgcolor=gray!5]{c}
static gpointer osc_reader_thread(gpointer data)
{
    AppState *st = data;
    uint8_t buf[8192];
    size_t have = 0;

    while (st->osc_thread_run) {
        ssize_t r = read(st->fd_osc, buf + have,
                         sizeof(buf) - have);
        if (r <= 0) {
            g_usleep(1000);
            continue;
        }
        have += (size_t)r;

        /* Поиск sync 0x55 0xAA */
        size_t pos = 0;
        while (have >= 8) {
            if (!(buf[pos] == 0x55 && buf[pos + 1] == 0xAA)) {
                pos++;
                have--;
                continue;
            }
            /* ... */
        }
    }
    return NULL;
}
\end{minted}
\caption{Поток чтения осциллографа}
\label{lst:arch-osc-reader}
\end{listing}

\textbf{Что здесь важно}:

\begin{itemize}
    \item \textbf{Отдельный поток}~--- \texttt{read()}
          \emph{блокируется}, пока нет данных. Если
          делать это в \emph{главном} потоке~---
          интерфейс \emph{зависнет}.
    \item \textbf{Буфер 8192 байт}~--- \emph{достаточно}
          для нескольких кадров \texttt{OSC\_DATA}.
    \item \textbf{Поиск sync}~--- \emph{синхронизация}
          по маркеру \texttt{0x55 0xAA}.
    \item \textbf{g\_usleep(1000)}~--- если данных
          нет~--- \emph{подождать} 1 мс, не \emph{жечь}
          CPU.
\end{itemize}

\begin{warnbox}
\textbf{Поток~--- это сложно.} Нужно \emph{синхронизировать}
доступ к общим данным через \texttt{GMutex}. Нужно
\emph{корректно} завершать поток. Нужно \emph{не} трогать
GTK из потока~--- только \texttt{gtk\_widget\_queue\_draw()}.
\textbf{Мы разберём это в главе 12.}
\end{warnbox}

\subsection{draw\_scope.c}

\textbf{Что делает}: \emph{рисует} осциллограмму через
\textbf{Cairo}.

\begin{itemize}
    \item \texttt{draw\_scope()}~--- \emph{колбэк}
          для \texttt{GtkDrawingArea}.
    \item \textbf{Масштабирование}~--- \emph{подгонка}
          данных под размер виджета.
    \item \textbf{Сетка}, \textbf{оси}, \textbf{метки}~---
          \emph{как} в \emph{первой} книге.
\end{itemize}

\subsection{ui.c}

\textbf{Что делает}: \emph{создаёт} весь интерфейс.

\begin{itemize}
    \item \texttt{app\_activate()}~--- \emph{главная}
          функция активации.
    \item \texttt{build\_scope\_tab()}~--- \emph{вкладка}
          осциллографа.
    \item \texttt{build\_generator\_tab()}~--- \emph{вкладка}
          генератора.
\end{itemize}

\begin{ideabox}
\textbf{Разделение ответственности~--- ключ.} \texttt{ui.c}
\emph{не знает}, как \emph{работает} порт. \texttt{serial\_port.c}
\emph{не знает}, как \emph{выглядит} интерфейс. Каждый
модуль~--- \emph{своя} задача.
\end{ideabox}

\subsection{main.c: короткий}

\textbf{Что делает}: \emph{точка} входа.

\begin{listing}[htbp]
\begin{minted}[frame=none, bgcolor=gray!5]{c}
int main(int argc, char **argv)
{
    AppState st = {0};
    g_mutex_init(&st.osc_lock);

    GtkApplication *app = gtk_application_new(
        "student.oscgen", G_APPLICATION_DEFAULT_FLAGS);
    g_signal_connect(app, "activate",
                     G_CALLBACK(app_activate), &st);

    int status = g_application_run(G_APPLICATION(app),
                                   argc, argv);
    g_object_unref(app);

    /* Остановить поток, закрыть порты */
    if (st.osc_thread) {
        st.osc_thread_run = false;
        g_thread_join(st.osc_thread);
    }
    if (st.fd_osc > 0) close(st.fd_osc);
    if (st.fd_gen > 0) close(st.fd_gen);

    return status;
}
\end{minted}
\caption{\texttt{main.c}~--- точка входа}
\label{lst:arch-main}
\end{listing}

\textbf{Что здесь важно}:

\begin{itemize}
    \item \texttt{gtk\_application\_new()}~--- \emph{новый}
          API GTK4. Про \emph{старый} (\texttt{gtk\_init})
          мы говорили в \emph{первой} книге.
    \item \texttt{g\_mutex\_init()}~--- \emph{инициализация}
          мьютекса \emph{до} запуска потока.
    \item \texttt{g\_thread\_join()}~--- \emph{корректное}
          завершение потока \emph{перед} выходом.
    \item \texttt{close(st.fd\_osc)}~--- \emph{закрытие}
          портов \emph{после} остановки потока.
\end{itemize}

\begin{keybox}
\textbf{Порядок завершения~--- важен.} Сначала~---
\emph{остановить} поток. Потом~--- \emph{закрыть} порты.
\emph{Наоборот}~--- поток может \emph{читать} из
\emph{закрытого} порта~--- и \emph{упадёт}.
\end{keybox}

%%%%%%%===============================





\section{Модули прошивки}
\label{sec:arch-firmware-modules}

ПК-приложение~--- это \emph{одна} программа. Прошивка~---
это \textbf{две} программы: для осциллографа и для генератора.
У каждой~--- \emph{свои} модули.

\subsection{Осциллограф (Nucleo-F401RE)}

Прошивка осциллографа состоит из \emph{четырёх} модулей.

\begin{table}[htbp]
\centering
\caption{Модули прошивки осциллографа}
\label{tab:arch-osc-modules}
\small
\begin{tabularx}{\textwidth}{@{}l L@{}}
\toprule
\textbf{Модуль} & \textbf{Ответственность} \\
\midrule
\texttt{adc.c}     & Настройка ADC: частота, разрешение, каналы \\
\texttt{dma.c}     & Двойной буфер (ping-pong), прерывания \\
\texttt{trigger.c} & Триггер: уровень, фронт, режим \\
\texttt{usb\_cdc.c} & Отправка кадров \texttt{OSC\_DATA}, приём команд \\
\bottomrule
\end{tabularx}
\end{table}

\textbf{Как это работает}:

\begin{enumerate}
    \item \texttt{adc.c} \emph{настраивает} ADC на \emph{частоту}
          дискретизации (например, \SI{100}{\kilo\hertz}).
    \item \texttt{dma.c} \emph{запускает} DMA~--- \emph{двойной}
          буфер (ping-pong). Пока \emph{одна} половина
          заполняется~--- \emph{другая} отдаётся прошивке.
    \item \texttt{trigger.c} \emph{ищет} \emph{начало} сигнала~---
          момент, когда напряжение \emph{пересекает} заданный
          уровень.
    \item \texttt{usb\_cdc.c} \emph{упаковывает} семплы в кадры
          \texttt{OSC\_DATA} и \emph{отправляет} их на ПК.
\end{enumerate}

\begin{keybox}
\textbf{Ключевая идея~--- ping-pong.} Пока DMA \emph{пишет}
в \emph{первую} половину буфера~--- прошивка \emph{читает}
\emph{вторую}. Когда первая \emph{заполнилась}~---
они \emph{меняются} местами. Так \emph{непрерывно} захватываем
сигнал, \emph{не} теряя данные.
\end{keybox}

\subsection{Генератор (Blue Pill)}

Прошивка генератора тоже состоит из \emph{четырёх} модулей.

\begin{table}[htbp]
\centering
\caption{Модули прошивки генератора}
\label{tab:arch-gen-modules}
\small
\begin{tabularx}{\textwidth}{@{}l L@{}}
\toprule
\textbf{Модуль} & \textbf{Ответственность} \\
\midrule
\texttt{dac.c}       & Настройка DAC: частота, разрешение \\
\texttt{pwm.c}       & Альтернатива DAC для меандра \\
\texttt{wave\_table.c} & Таблицы форм (синус, пила, треугольник) \\
\texttt{usb\_cdc.c}  & Приём команд, отправка ответов \\
\bottomrule
\end{tabularx}
\end{table}

\textbf{Как это работает}:

\begin{enumerate}
    \item \texttt{usb\_cdc.c} \emph{принимает} команды
          (set\_wave, set\_freq, set\_ampl).
    \item \texttt{wave\_table.c} \emph{заполняет} таблицу
          формы (256 или 1024 точки).
    \item \texttt{dac.c} \emph{запускает} DMA circular~---
          таблица \emph{воспроизводится} \emph{бесконечно}.
    \item \texttt{pwm.c} \emph{используется} для \emph{меандра}~---
          когда нужна \emph{точная} частота и \emph{скважность}.
\end{enumerate}

\begin{ideabox}
\textbf{Почему \emph{два} способа~--- DAC и PWM?} DAC даёт
\emph{плавный} сигнал (синус, пила). PWM~--- \emph{прямоугольник}
с \emph{регулируемой} скважностью. Для \emph{меандра} PWM
\emph{лучше}~--- он \emph{точнее} и \emph{проще}.
\end{ideabox}

\section{Границы проекта}
\label{sec:arch-boundaries}

\textbf{Важно} \emph{обозначить}, что \emph{входит} в проект,
а что~--- \emph{нет}. Иначе читатель \emph{ждёт} большего,
чем мы \emph{можем} дать.

\subsection{Что мы \emph{делаем}}

\begin{itemize}
    \item \textbf{Осциллограф} на \SI{500}{\kilo\hertz}~---
          \emph{один} канал, \SI{12}{\bit}.
    \item \textbf{Триггер} по уровню и фронту~---
          \emph{простой}, но \emph{рабочий}.
    \item \textbf{Генератор} до \SI{100}{\kilo\hertz}~---
          \emph{шесть} форм.
    \item \textbf{Протокол}~--- \emph{бинарный}, с CRC16.
    \item \textbf{ПК-приложение} на GTK4~--- \emph{потоки},
          \emph{Cairo}, \emph{управление}.
\end{itemize}

\subsection{Что мы \emph{не} делаем}

\begin{itemize}
    \item \textbf{Не} пишем \emph{полноценный} осциллограф
          на \SI{150}{\mega\hertz} (как HANTEK).
    \item \textbf{Не} делаем \emph{многоканальный} осциллограф
          (\emph{только} один канал).
    \item \textbf{Не} поддерживаем \emph{все} формы сигналов
          (\emph{шесть}~--- достаточно).
    \item \textbf{Не} делаем \emph{FFT} (спектральный анализ).
    \item \textbf{Не} пишем \emph{драйверы} ядра~---
          \emph{используем} стандартные.
    \item \textbf{Не} поддерживаем \emph{Windows}~---
          \emph{только} Linux.
\end{itemize}

\begin{warnbox}
\textbf{Честность~--- важна.} Мы \emph{не} обещаем
<<профессиональный осциллограф>>. Мы обещаем
\textbf{учебный проект}, который \emph{работает} и
\emph{объясняет}, как устроены \emph{настоящие} приборы.
\end{warnbox}

\section{Уроки главы}
\label{sec:arch-lessons}

\begin{enumerate}
    \item \textbf{Три мира:} осциллограф, генератор, ПК.
          У каждого~--- \emph{своя} задача.
    \item \textbf{Потоки данных:} команды~--- вниз,
          данные~--- вверх.
    \item \textbf{Монолитный \texttt{main.c}~--- плохо.}
          \emph{Шесть} модулей на ПК~--- лучше.
    \item \textbf{Поток чтения~--- сложно, но необходимо.}
          Без него интерфейс \emph{зависнет}.
    \item \textbf{Разделение ответственности~--- ключ.}
          Каждый модуль делает \emph{одну} вещь.
    \item \textbf{Ping-pong DMA}~--- \emph{непрерывный}
          захват сигнала.
    \item \textbf{Границы проекта~--- важно обозначить.}
          Честность с читателем.
\end{enumerate}

\section{Чек-лист грабель}
\label{sec:arch-pitfalls}

\begin{table}[htbp]
\centering
\caption{Грабли архитектуры}
\label{tab:arch-pitfalls}
\small
\begin{tabularx}{\textwidth}{@{}c L L@{}}
\toprule
\textbf{№} & \textbf{Грабля} & \textbf{Решение} \\
\midrule
1 & Монолитный \texttt{main.c} & Разбить на модули \\
2 & Чтение в главном потоке & Отдельный поток \\
3 & Гонка данных & \texttt{GMutex} \\
4 & Зависание при закрытии & \texttt{g\_thread\_join()} \\
5 & Забыли \texttt{g\_mutex\_init()} & Инициализация до потока \\
6 & Закрытие порта до остановки потока & Порядок важен \\
7 & Модуль делает слишком много & Разделить \\
8 & Нет границ проекта & Обозначить \\
\bottomrule
\end{tabularx}
\end{table}

\section{Упражнения}
\label{sec:arch-exercises}

\begin{exercisebox}
\textbf{Упражнение 1.} Нарисуйте \emph{диаграмму} своих
модулей для \emph{своего} проекта. Какие модули у вас есть?
Как они связаны?
\end{exercisebox}

\begin{exercisebox}
\textbf{Упражнение 2.} Разбейте \emph{монолитный}
\texttt{main.c} из репозитория на \emph{модули}.
Сколько файлов получилось? Какие функции в каждом?
\end{exercisebox}

\needspace{5\baselineskip}
\begin{exercisebox}
\textbf{Упражнение 3.} Объясните \emph{одним предложением},
за что отвечает каждый модуль \emph{вашего} проекта.
Если не получается~--- \emph{пересмотрите} границы.
\end{exercisebox}

\begin{exercisebox}
\textbf{Упражнение 4.} Предложите \emph{свои} границы проекта~---
что бы вы \emph{добавили} в наш? Что \emph{убрали}?
\end{exercisebox}

\section{Проверьте сами}
\label{sec:arch-check}

\begin{enumerate}
    \item Откройте репозиторий \texttt{dualboard-stm32-oscilloscope}.
    \item Найдите \texttt{main.c} для ПК.
    \item Посчитайте, сколько \emph{функций} в нём.
    \item Разбейте их на \emph{шесть} модулей (список).
    \item Нарисуйте \emph{диаграмму} зависимостей.
\end{enumerate}

\begin{ideabox}
\textbf{Если получилось~--- вы поняли архитектуру.} Это
\emph{главное}. Код~--- \emph{следствие}. Архитектура~---
\emph{причина}.
\end{ideabox}




